% References -- consolidated thesis-wide bibliography (O-RAN / BMPP-DQN
% track, Chapters 1-5).
%
% This replaces the per-chapter local \begin{thebibliography} blocks each
% chapter file previously carried (each chapter's own header comment
% flagged this merge as pending until a thesis-wide main.tex existed). All
% \bibitem keys below are verbatim unions of what each chapter cited;
% every \cite{} call across Chapters 1-5 resolves here. Do not duplicate a
% key back into an individual chapter file -- that is exactly the
% drift/duplicate-label risk this merge resolves (previously surfaced as
% LaTeX's own "multiply defined" label warnings when all chapters were
% compiled together).
%
% \renewcommand{\bibname}{References} overrides the thebibliography
% environment's default report-class heading ("Bibliography") to match
% this file's own title below and this track's citation convention
% (IEEE numbered style).

\renewcommand{\bibname}{References}

\begin{thebibliography}{99}

\bibitem{qazzaz2026oreo}
M.~M.~H. Qazzaz, A. Salama, M. Hafeez, and S.~A.~R. Zaidi,
``OREO: Open RAN Energy Optimization via Deep Reinforcement Learning for 6G
Networks,'' \emph{IEEE Open Journal of the Communications Society}, vol.~7,
pp.~4165--4182, 2026.

\bibitem{oran2021mvp}
O-RAN Alliance, ``O-RAN Minimum Viable Plan and Commercialization
Strategy,'' Jun. 2021.

\bibitem{xiong2018pdqn}
J. Xiong, Q. Wang, Z. Yang, P. Sun, L. Han, Y. Zheng, H. Fu, T. Zhang, J. Liu,
and H. Liu, ``Parametrized Deep Q-Networks Learning: Reinforcement Learning
with Discrete-Continuous Hybrid Action Space,'' arXiv preprint
arXiv:1810.06394, 2018.

\bibitem{bester2019mpdqn}
C.~J. Bester, S.~D. James, and G.~D. Konidaris, ``Multi-Pass Q-Networks for
Deep Reinforcement Learning with Parameterised Action Spaces,'' arXiv
preprint arXiv:1905.04388, 2019.

\bibitem{tavakoli2018branching}
A. Tavakoli, F. Pardo, and P. Kormushev, ``Action Branching Architectures for
Deep Reinforcement Learning,'' in \emph{Proc. AAAI Conference on Artificial
Intelligence}, 2018.

\bibitem{bordin2025design}
M. Bordin, A. Lacava, M. Polese, S. Satish, M. AnanthaSwamy Nittoor,
R. Sivaraj, F. Cuomo, and T. Melodia, ``Design and Evaluation of Deep
Reinforcement Learning for Energy Saving in Open RAN,'' in \emph{Proc. IEEE
22nd Consumer Communications \& Networking Conference (CCNC)}, Las Vegas,
NV, USA, Jan. 2025, pp.~1--6.

\bibitem{sthankiya2024survey}
A. Sthankiya et al., ``A Survey on AI-Driven Energy Optimization in
Next-Generation RAN,'' arXiv preprint arXiv:2411.02164, 2024.

\bibitem{liang2026federated}
X. Liang, A. Al-Tahmeesschi, S.~B. Chetty, and H. Ahmadi, ``Green O-RAN
Operation: a Modern ML-Driven Network Energy Consumption Optimization,'' in
\emph{Proc. IEEE GLOBECOM}, 2025, pp.~4475--4480.

\bibitem{sohaib2024transfer}
M. Sohaib et al., ``DRL Transfer Learning for Cloud-Native O-RAN,''
arXiv preprint arXiv:2407.11563, 2024.

\bibitem{iqbal2021ddqn}
A. Iqbal, M.-L. Tham, and Y.~C. Chang, ``Double Deep Q-Network-Based
Energy-Efficient Resource Allocation in Cloud Radio Access Network,''
\emph{IEEE Access}, vol.~9, pp.~20440--[end page not independently
confirmed in this repository], 2021, doi: 10.1109/ACCESS.2021.3054909.

\bibitem{fathy2021ann}
M. Fathy, M.~S. Abood, and M.~M. Hamdi, ``Optimization of Energy-Efficient
Cloud Radio Access Networks for 5G using Neural Networks,'' in \emph{Proc.
Int. Conf. Intelligent Technology, System and Service for Internet of
Everything (ITSS-IoE)}, 2021.

\bibitem{alzubaedi2019cran}
W.~H.~A. Al-Zubaedi, ``Planning a C-RAN Deployment for the Next Generation
Cellular Networks,'' Ph.D. dissertation, Dept.\ of Electronic and Computer
Engineering, Brunel University London, London, U.K., 2019.

\bibitem{shengren2022benchmark}
H. Shengren, E.~M. Salazar Duque, P.~P. Vergara, and P. Palensky,
``Performance Comparison of Deep RL Algorithms for Energy Systems Optimal
Scheduling,'' in \emph{Proc. IEEE PES Innovative Smart Grid Technologies
Conference Europe (ISGT-Europe)}, 2022, pp.~1--6.

\bibitem{ieee2024selfoptimized}
``Self-Optimized Agent for Load Balancing and Energy Efficiency: A
Reinforcement Learning Framework With Hybrid Action Space,'' \emph{IEEE
Trans.\ on Network and Service Management}, vol.~21, no.~4, pp.~4902--4919,
Jul. 2024.

\bibitem{sharma2023massivemimo}
S. Sharma and W. Yoon, ``Energy Efficient Power Allocation in Massive MIMO
Based on Parameterized Deep DQN,'' \emph{Electronics}, vol.~12, no.~21,
p.~4517, 2023.

\bibitem{kabir2023mobility}
H. Kabir, M.-L. Tham, and Y.~C. Chang, ``Mobility-Aware Resource Allocation
in IoT Network for Post-Disaster Communications with Parameterized
Reinforcement Learning,'' \emph{Sensors}, vol.~23, no.~14, p.~6448,
Jul. 2023.

\bibitem{luyanzeng2026eexapp}
J. Lu, P. Yan, and H. Zeng, ``EExApp: GNN-Based Reinforcement Learning for
Radio Unit Energy Optimization in 5G O-RAN,'' in \emph{Proc.\ IEEE
INFOCOM}, 2026, pp.~1--10.

\bibitem{yan2025fluctuates}
C. Yan et al., ``Hybrid Reinforcement Learning in Parameterized Action
Space via Fluctuates Constraint,'' \emph{Engineering Applications of
Artificial Intelligence}, vol.~162, p.~112499, 2025.

\bibitem{rezazadeh2022actorcritic}
F. Rezazadeh, H. Chergui, L. Christofi, and C. Verikoukis, ``Actor-Critic-
Based Learning for Zero-Touch Joint Resource and Energy Control in Network
Slicing,'' arXiv preprint, 2022.

\bibitem{amiri2022vnf}
E. Amiri, N. Wang, M. Shojafar, and R. Tafazolli, ``Energy-Aware Dynamic
VNF Splitting in O-RAN Using Deep Reinforcement Learning,'' in \emph{Proc.
IEEE Int.\ Conf.\ Local Computer Networks (LCN)}, 2022, pp.~422--429.

\bibitem{barker2025mec}
R. Barker, T. Seyfi, and F. Afghah, ``Advancements in Mobile Edge Computing
and Open RAN: Leveraging Artificial Intelligence and Machine Learning for
Wireless Systems,'' \emph{arXivLens}, 2025.

\bibitem{lassoued2026review}
N. Lassoued and N. Boujnah, ``A Comprehensive Review of Energy Efficiency in
5G Networks: Past Strategies, Present Advances, and Future Research
Directions,'' \emph{Computers}, vol.~15, no.~1, p.~50, 2026.

\bibitem{chuang2025hybrid}
Chuang et al., ``[Title and venue not independently verified in this
repository -- known only via a descriptive secondary reference, per
AGENTS.md's Foundational References table:] Hybrid A3C and Dueling DQN
applied to industrial IoT task scheduling in a 5G Cloud-RAN setting,''
2025. \emph{Needs primary-source verification before this entry is used
outside this draft.}

\bibitem{threegpp38801}
3GPP, ``Study on New Radio Access Technology: Radio Access Architecture and
Interfaces,'' TR 38.801, V0.4.0, Aug. 2016.

\bibitem{threegpp38864}
3GPP, ``Study on Network Energy Savings for NR,'' TR 38.864, 2024.

\bibitem{auer2011earth}
G. Auer et al., ``How Much Energy Is Needed to Run a Wireless Network?,''
\emph{IEEE Wireless Communications}, vol.~18, no.~5, pp.~40--49, Oct. 2011.

\bibitem{benetel2024ran550}
Benetel, ``RAN550 Outdoor/Indoor O-RU Datasheet (n78/n79, Split 7.2x),''
product datasheet. Max.\ TX output power (total EIRP): 2\,W (33\,dBm).

\bibitem{hwang1981topsis}
C.-L. Hwang and K. Yoon, \emph{Multiple Attribute Decision Making: Methods
and Applications}. Berlin, Germany: Springer-Verlag, 1981.

\bibitem{zeleny1982mcdm}
M. Zeleny, \emph{Multiple Criteria Decision Making}. New York, NY, USA:
McGraw-Hill, 1982.

\bibitem{opricovic2004vikor}
S. Opricovic and G.-H. Tzeng, ``Compromise solution by MCDM methods: A
comparative analysis of VIKOR and TOPSIS,'' \emph{European Journal of
Operational Research}, vol.~156, no.~2, pp.~445--455, 2004.

\bibitem{cohen1988statistical}
J. Cohen, \emph{Statistical Power Analysis for the Behavioral Sciences},
2nd ed. Hillsdale, NJ, USA: Lawrence Erlbaum Associates, 1988.

\bibitem{li2022hyar}
Y. Li, Q. Zhang, Y. Wang, H. Jianye, and Y. Yu, ``HyAR: Addressing
Discrete-Continuous Action Space via Hybrid Action Representation,'' in
\emph{Proc. International Conference on Learning Representations (ICLR)},
2022.

\end{thebibliography}
